\section{The Tessellation Catalog}
\label{app:catalog}

The regular hyperbolic honeycombs are a finite, fully enumerated family, and the substrate is one entry on it. The enumeration is governed by the same angle condition that fixes any Coxeter tessellation, and it thins out with dimension. The compact regular honeycombs, those with finite cells and finite vertex figures, exist in two, three, and four dimensions and stop there. The paracompact ones, with ideal cells or vertices, reach into the fourth and fifth dimensions and stop at the fifth. There is nothing regular beyond the fifth dimension at all, which is one half of why the substrate is four-dimensional.

The survey of the body measured every buildable member the same way, and the two summary numbers are worth recording here. A propagating Kahler-Dirac fermion lives on all forty-two buildable tessellations, with a clean return probability falling with size on each, so matter is a universal feature of negative curvature. The native spinor sites, the directions that carry spin without a detour, appear on exactly seven, the honeycombs faceted by the $24$-cell, and these are dimension-gated, first appearing in four dimensions. The members that matter most for the paper are few.

\begin{table}[ht]
\centering
\small
\begin{tabular}{@{}llp{0.40\textwidth}@{}}
\toprule
honeycomb & dimension & role \\
\midrule
$\{7,3\}$ & 2 & the flat heptagrid face, drawn for intuition, the cleanest holography \\
$\{5,3,4\}$ & 3 & the comparison substrate, three-dimensional cusp, no spinor in its sites \\
$\{3,4,3,4\}$ & 4 & the committed substrate, crystallographic, spinor-carrying, flat cubic cusp \\
$\{4,3,4,3\}$ & 4 & its dual, the same sites, equally crystallographic \\
$\{3,4,3,3\}$ & 4 & the flat twin, the same spinor sites but no curvature, so excluded \\
$\{3,4,3,3,4\}$ & 5 & the pentacomb, spinor sites and curvature together, resolving the trade \\
\bottomrule
\end{tabular}
\caption{The members of the family that figure in the paper. The full per-tessellation measurements, degree, spectral dimension, crystallography, and spinor content, are produced by the simulator and kept with the code.}
\label{tab:catalog}
\end{table}

The committed substrate is the four-dimensional member that keeps the spinor sites of the $24$-cell while delivering a flat three-dimensional cusp, and the selection argument of the body shows it to be the only one that meets all four demands at once. The catalog's role is to make that uniqueness concrete, by laying out the small finite field within which the choice is forced.
